% Generated draft: requires mathematical and source review.
\documentclass[11pt]{article}
\usepackage[margin=25mm]{geometry}
\usepackage{fontspec}
\setmainfont{Latin Modern Roman}
\usepackage{amsmath,amssymb,mathtools,unicode-math}
\setmathfont{Latin Modern Math}
\usepackage{enumitem,needspace,xurl,hyperref}
\setlength{\parindent}{0pt}
\setlength{\parskip}{5pt}
\setlength{\emergencystretch}{3em}
\newcommand{\N}{\mathbb{N}}
\newcommand{\Z}{\mathbb{Z}}
\newcommand{\Q}{\mathbb{Q}}
\newcommand{\R}{\mathbb{R}}
\newcommand{\C}{\mathbb{C}}
\newcommand{\F}{\mathbb{F}}
\newcommand{\A}{\mathbb{A}}
\newcommand{\PP}{\mathbb{P}}
\newcommand{\E}{\mathbb{E}}
\newcommand{\eps}{\varepsilon}
\newcommand{\dd}{\,\mathrm{d}}
\begin{document}
\section*{UnsolvedMath ID: 701063}
\textbf{Source problem code:} catalogue entry
\subsection{Definitions and mathematical statement}
All groups considered are finite. A formation is a nonempty class of groups, closed under isomorphisms, homomorphic images, and finite subdirect products. Explicitly, the last condition means that if \(N_1,\ldots,N_r\trianglelefteq G\), with \(r\geq 1\), satisfy \(\bigcap_{i=1}^r N_i=1\) and \(G/N_i\in\mathfrak F\) for every \(i\), then \(G\in\mathfrak F\).

Use the following standard meaning of a local formation. A local definition assigns to each prime \(p\) a formation \(f(p)\), with the empty class also permitted as a value. The class \(\operatorname{LF}(f)\) consists of the finite groups \(G\) such that, for every chief factor \(A/B\) of \(G\) and every prime \(p\mid |A/B|\),
\[
G/C_G(A/B)\in f(p).
\]
Here a chief factor has \(B<A\), both normal in \(G\), with no normal subgroup of \(G\) strictly between them, and
\[
C_G(A/B)=\{g\in G: gag^{-1}B=aB\text{ for every }a\in A\}.
\]
A formation is local if it equals \(\operatorname{LF}(f)\) for some such assignment \(f\).

For formations \(\mathfrak H\subseteq\mathfrak M\), write
\[
[\mathfrak H,\mathfrak M]
=\{\mathfrak F:\mathfrak F\text{ is a formation and }\mathfrak H\subseteq\mathfrak F\subseteq\mathfrak M\}.
\]
This lattice is ordered by inclusion. Its meet is intersection, and \(\mathfrak F_1\vee\mathfrak F_2\) is the smallest formation containing \(\mathfrak F_1\cup\mathfrak F_2\).

\textbf{Question (A. N. Skiba).} For every pair of local formations \(\mathfrak H\subseteq\mathfrak M\) and every pair of formations \(\mathfrak F_1,\mathfrak F_2\in[\mathfrak H,\mathfrak M]\) satisfying
\[
\mathfrak F_1\cap\mathfrak F_2=\mathfrak H,
\qquad
\mathfrak F_1\vee\mathfrak F_2=\mathfrak M,
\]
must both \(\mathfrak F_1\) and \(\mathfrak F_2\) be local formations? All inclusions are non-strict; endpoint and degenerate cases are included.

The supplied source states that the answer is affirmative when \(\mathfrak H=(1)\), where \((1)\) denotes the formation consisting of the trivial group. This special-case assertion is retained as supplied and is not independently verified here.
\subsection{Short English statement}
If two formations complement each other in an interval whose least and greatest formations are local, must both complementary formations also be local? The supplied source states that this holds when the least formation consists only of the trivial group.
\subsection{Sources}
\begin{itemize}
\item Local UnsolvedMath registry, permanent ID \(701063\), source problem code \texttt{catalogue entry}, title \emph{Kourovka Notebook Problem 14.78}; supplied complete source record, upstream revision \texttt{d371ee917ac9647fbb06cd6fc529deec6f3dda82}.
\item \emph{The Kourovka Notebook (Unsolved Problems in Group Theory)}, Problem 14.78, attributed to A. N. Skiba, proposed in 1999 in the 14th issue. The supplied background also identifies the 21st issue edition (2026), \texttt{arXiv:1401.0300v48}. The supplied collection description names E. I. Khukhro and V. D. Mazurov as editors.
\item The registry labels this problem \emph{open}. This is a supplied status label, not an independently verified assessment. The supplied statement also records an affirmative answer for \(\mathfrak H=(1)\), without a separate reference for that special case.
\end{itemize}
\end{document}
